\section{Introduction}\label{sec:intro}

The finite element method (FEM) is the standard tool for structural analysis \cite{zienkiewicz2013fem,bathe2014fep}.
Its cost grows with the number of degrees of freedom (DOF), and many-query tasks such as parameter studies, design
optimisation, uncertainty quantification and real-time prediction need the same model solved many times.
Model order reduction (MOR) addresses this by replacing the full model with a small one built from a projection basis
\cite{benner2015survey,hesthaven2016rb,antoulas2005}.
In practice the two steps live in different software: one tool assembles and solves the full model, another builds the
reduced model, and the user moves matrices between them.

This paper describes \texttt{computation-suite}, an open-source collection of two Python packages that keep both steps
in one code base while keeping them decoupled:

\begin{itemize}[leftmargin=*,itemsep=2pt]
  \item \fea{} (version~\feaVersion{}) is a finite element package for structural mechanics. It registers \nElements{}
  element formulations and \nConstitutive{} constitutive models and provides static, modal, harmonic, transient and
  nonlinear analyses (\cref{sec:arch}).
  \item \rom{} (version~\romVersion{}) works on plain arrays $\K,\M,\C,\F$. It includes projection-based, state-space,
  surrogate-based and hyper-reduction methods and component mode synthesis, and can be used on matrices that come from any
  other source.
\end{itemize}

Both depend only on NumPy \cite{harris2020numpy} and SciPy \cite{virtanen2020scipy}. PyTorch \cite{paszke2019pytorch} is an optional
backend for GPU solves and differentiable stress recovery.
The library code has \totalLines{} lines (\feaLines{} in \fea{}, \romLines{} in \rom{}) and is checked by \feaTests{} and
\romTests{} automated tests.

\paragraph{Contributions.}
The paper (i) documents the architecture and the interfaces between the FE and ROM layers;
(ii) states the mathematics and algorithms behind the main components, with references to the original methods;
(iii) reports verification against closed-form and independent reference solutions, including observed convergence orders;
(iv) reports speed-ups of reduced and batched computations measured by scripts that ship with the repository, so the numbers
can be reproduced; and (v) states the known limitations. We make no claim that the package is faster than, or replaces,
established libraries; \cref{sec:related} places it among them.

\paragraph{Outline.}
\Cref{sec:related} reviews related software and \cref{sec:arch} describes the package structure and interfaces. \Cref{sec:fe,sec:constraints,sec:analysis,sec:nonlinear,sec:recovery} give the finite element formulation, constraints, linear and nonlinear analyses and derived results;
\cref{sec:rom} gives the reduction methods. \Cref{sec:verif,sec:perf} report verification and performance, \cref{sec:usage} shows a workflow, and \cref{sec:limits,sec:avail} state limitations and availability.
